\chapter{Dashboard Streamlit}
\label{chap:dashboard}

\section{Architecture}

Le livrable de restitution est \filename{dashboard.py}, lancé par \texttt{streamlit run dashboard.py}. Il s'appuie sur Streamlit, Pandas et Matplotlib uniquement \cite{streamlitdocs}. Le DataFrame source est chargé une fois depuis PostgreSQL (\texttt{@st.cache\_data}), en lecture seule. Les filtres produisent un DataFrame distinct ; le DataFrame chargé n'est jamais modifié.

\section{Filtres et KPI}

La barre latérale propose les filtres Company, Campaign Type, Channel, Target Audience, Customer Segment, Location, Quarter et Month, chacun avec l'option « Tous ». Les filtres se combinent. À chaque changement, les KPI et les graphiques concernés sont recalculés sur le sous-ensemble.

Les cartes d'indicateurs affichent : nombre de campagnes, ROI moyen, coût d'acquisition moyen, taux de conversion moyen (en pourcentage), CTR moyen (en pourcentage, par exemple 14,04~\% et non 0,1404) et engagement moyen.

\section{Organisation des vues}

Le dashboard reprend la logique des visualisations validées, recalculée dynamiquement :

\begin{enumerate}
\item vue d'ensemble (volumes entreprise, type, canal) ;
\item distributions (ROI, coût, conversion, CTR) ;
\item audience, segments et canaux (heatmaps ROI et classes de CTR) ;
\item analyse temporelle (volume mensuel et ROI mensuel, avec échelle prudente) ;
\item corrélations et analyses complémentaires (durée, engagement) ;
\item constats et limites issus de l'EDA.
\end{enumerate}

Un export CSV optionnel, nommé \filename{marketing\_campaigns\_filtered.csv}, est généré uniquement à partir du DataFrame filtré.

\section{Emplacement pour une capture}

Aucune capture d'écran du dashboard n'est jointe à ce rapport. L'encadré suivant réserve un emplacement pour une future illustration.

\begin{figure}[H]
\centering
\fbox{%
\begin{minipage}{0.88\textwidth}
\centering
\vspace{1.6cm}
\textit{Emplacement réservé --- capture d'écran du dashboard Streamlit}\\[0.3cm]
{\small Commande : \texttt{streamlit run dashboard.py}}
\vspace{1.6cm}
\end{minipage}}
\caption{Emplacement prévu pour une capture du dashboard (non disponible au moment de la rédaction)}
\label{fig:dashboard-placeholder}
\end{figure}
